% Pursuing Stacks (cotes 134-2 to 134-8): archive cote 134-6, site page 20 (PDF page 21), typescript page 339.
% typescript. Draft compared: ps5.tex:2376-2437.
% AI-transcribed by Claude Opus 5.5 (claude-opus-5-5, adaptive thinking, effort high), Message Batches API, batch ps3 (msgbatch_01XdFmuVVBAcgc7EFxdtCqsj), 2026-11-12; not yet checked by eye.
\scanpage{134-6/20}
\mspage{339}

cohomology invariants.  This is the question \over{\unsure{or|on}}{of} \uline{how far a homotopy
type can be expressed in terms of homology or cohomology invariants}
(or both together),  \uuline{plus} \uline{some relevant extra structure},  the most im-%
portant surely being cup-products in cohomology (or, dually, ``interior''
operation of cohomology on homology).  Once the notion of derived catego-%
ries of various kinds had become familiar,  in the early sixties,  the
question would appear as expressing\hand{,}or recovering\hand{,}a homotopy type,  name-%
ly an object in the (highly non abelian) ``derived category'' (Hot),  in
terms of it's abelianization in  Hotab $=$ D.(Ab),  \uline{endowed with suitable
extra structure} .  It was about clear that this extra structure had to
include,  as it's main non commutative item,  the fundamental group\inserted{$\pi$},
so as to allow for description of homology and cohomology invariants
with twisted coefficients.  The most natural candidate for expressing
this would be the chain complex associated to the universal covering,
viewed as an object in the derived category
\[ \text{(26)} \qquad \mathrm{D.}\;(\;\underline{Z}\;(\hand{\pi})\;) \]
of chain complexes of modules over the group ring $\underline{Z}\hand{(\pi)}$.  Another impor-%
tant structural item,  giving rise to all cup-products with non twisted
coefficients,  is the diagonal map  for the abelianization
\[ \mathrm{L.}\quad\text{in}\quad \mathrm{Hotab} = \mathrm{D.(Ab)} \]
namely a map
\[ \text{(27)} \qquad \mathrm{L.} \hand{\longrightarrow} \mathrm{L.} \overset{L}{\otimes} \mathrm{L.} \qquad , \]
where $\overset{L}{\otimes}$ is the\hand{``}total\hand{''}left derived functor of tensor product.  This
map is subjected to \xout{a} suitable  conditions,  concerning mainly commuta-%
tivity and associativity.  In case of a non \xout{simply} 1-connected space
i.e. $\hand{\pi} \neq 1$ ,  it should'nt be hard combining the two structural items
so as to get a structure embodying at any rate cup-products with arbitra-%
ry twisted coefficients.   One key question then in my mind,  which I
never really looked into,  was whether these two structures were enough
in order to reconstruct entirely (up to canonical isomorphism) the
(pointed, 0-connected)  homotopy type  giving rise to it ,  and hence
also any other homotopy invariants,  such as ``operations'' on cohomology
and the like,  K-invariants etc. \hand{$\|$}\margin{$\ulcorner$}If I got it right,  it \xout{is} has been
known now for quite a while that even for a 1-connected homotopy type,
so that the relevant structure reduces to (27),  \inserted{that} this is \uline{not} quite
enough for recovering the homotopy type,  maybe not even the rational
homotopy type.  I believe I first got this from Sullivan,  namely that

\margin{20}
